% 1. Настройки абзацного отступа
\setlength{\parindent}{1.25cm}

% 2. Глобальная настройка всех списков (enumitem)
\setlist{
	leftmargin=0pt,
	nosep,
	wide=\parindent,
	labelsep=0.5em,
	listparindent=\parindent
}
\setlist[itemize,1]{label=--} 
\setlist[enumerate,1]{label=\arabic*.}

% 3. НАСТРОЙКА РАЗМЕРА И ШРИФТА ЗАГОЛОВКОВ (ФИКС ОГРОМНОГО ТЕКСТА)
% Устанавливаем normalsize (14pt) для всех типов заголовков
\setkomafont{disposition}{\rmfamily\bfseries\normalsize}
\setkomafont{chapter}{\rmfamily\bfseries\normalsize}
\setkomafont{section}{\rmfamily\bfseries\normalsize}
\setkomafont{subsection}{\rmfamily\bfseries\normalsize}

% Настройка оглавления (чтобы слово СОДЕРЖАНИЕ не было огромным)
\setkomafont{chapterentry}{\rmfamily\normalsize}

% 4. Настройка отступов и интервалов заголовков через KOMA-Script
% indent=1.25cm вернет "красную строку" заголовкам
\RedeclareSectionCommand[
    beforeskip=12pt, 
    afterskip=6pt, 
    indent=1.25cm
]{chapter}

\RedeclareSectionCommand[
    beforeskip=12pt, 
    afterskip=6pt, 
    indent=1.25cm
]{section}

\RedeclareSectionCommand[
    beforeskip=12pt, 
    afterskip=6pt, 
    indent=1.25cm
]{subsection}

% 5. Команда против переносов
\newcommand{\nohyphens}{\hyphenpenalty=10000\exhyphenpenalty=10000\sloppy}

% 6. Настройка подписей (caption)
\captionsetup{
	justification=centering,
	singlelinecheck=false
}
\DeclareCaptionLabelFormat{simple}{#1~#2}
\captionsetup{labelformat=simple}

% 7. Настройка таблиц
\DeclareCaptionLabelSeparator{emdash}{ --- }
\captionsetup[table]{
	labelsep=emdash,
	justification=raggedright,
	singlelinecheck=false,
	position=top,
	skip=6pt
}

% 8. Стили для листингов (JSON и TypeScript)
\lstset{
    basicstyle=\ttfamily\footnotesize,
    breaklines=true,
    texcl=true % Важно для поддержки кириллицы в комментариях кода
}

\lstdefinelanguage{json}{
	basicstyle=\ttfamily\footnotesize,
	numbers=none,
	stringstyle=\color{blue},
}

% Настройка библиографии под СФУ
\DefineBibliographyStrings{russian}{
	urlseen = {дата обращения:},
}